\scanpage{16}
\mspage{28}
\deleted{
\uline{Lemme} 5.12. Soient $S$ un préschéma, $R$ un préschéma constant tordu sur $S$.
Alors tout sous-préschéma fermé \deleted{$\unsure{T|Z}$}\inserted{$Z$} de $R$ qui est
\deleted{de type fini}\inserted{quasi-compact} sur $S$,
est fini sur $S$.

En effet, on est ramené au cas où $R$ est constant, donc de la forme
$I_S$, où $I$ est un ensemble, donc réunion filtrante croissante des $J_S$, où
$J$ parcourt les parties finies de $I$. On peut supposer de plus $S$ affine,
alors \deleted{$\unsure{R|Z}$} $Z$ est quasi-compact, donc contenu dans un des $J_S$. Comme $J_S$
est fini sur $S$, il en est de même de $Z$.

Le lemme 5.12. établit déjà l'assertion entre parenthèses dans
5.11. (ii). \deleted{L'équivalence $(ii)\unsure{\Leftrightarrow|-}(iii)$ est}\inserted{L'implication $(ii)\Rightarrow(iii)$, et \unsure{llimplication|l'implication} inverse}
dans le cas $S$ localement
noethérien est alors claire. Prouvons que $(i)\Rightarrow(ii)$, pour ceci, soit
$S'\longrightarrow S$ un morphisme étale fini surjectif qui splitte $H$ donc $R$, de sorte
que $R'$ s'écrit \deleted{comme} $M_S=\deleted{\unsure{\sum}}\coprod_{m\in M}R'_m$, où \deleted{tous} les $R'_m$ sont des ouverts
disjoints de $R'$, \inserted{$S'$-}isomorphes à $S'$. Soit $g:R'\longrightarrow R$ la projection, qui est
un morphisme fini étale surjectif, \inserted{donc un morphisme ouvert et fermé;}
alors les $R_m=g(R'_m)$ sont des parties
à la fois ouvertes et fermées de $R$, et évidemment quasi-compacts sur $S$
puisque les $R_m$ le sont. \deleted{Enfin, prouvons $(ii)\unsure{-|\Rightarrow}(i)$}

Enfin, prouvons $(ii)\Rightarrow(i)$ lorsque $H$ est de type fini sur $S$,\inserted{$S$ quasi compact ou connexe.} Soit
$M$ son type, on peut écrire $M=\deleted{\unsure{N}}\underline{Z}^{,r}\times N$, où $r$ est un entier $\geq0$
et $N$ un groupe fini. Considérons le préschéma
\[
P=\underline{\operatorname{Isom}}_{S\text{-gr}}(H,G)\subset
\underline{\operatorname{Hom}}_{S\text{-gr}}(H,G)\unsure{=|\ne}Q
\]
où $G=D_S(M)$, cf 5.8. et 5.10. \deleted{Alors $Q$ est isomorphe} On a un isomorphis\unsure{me}
\[
Q\simeq\underline{\operatorname{Hom}}_{S\text{-gr}}(G,\underline{G}_m^{,r})\times
\underline{\operatorname{Hom}}_{S\text{-gr}}(G,N)\simeq R^r\times E,
\]
où $E=\underline{\operatorname{Hom}}_{S\text{-gr}}(G,N)$ est fini sur $S$. Il s'ensuit que $Q$ est
\deleted{réunion de}\inserted{réunion de sous-}préschémas à la fois ouverts et fermés $Q_i$ finis sur $S$. \deleted{De même, il
s'ensuit que}\inserted{Donc} $P$ est réunion \deleted{\unsure{d'un}} de \deleted{\unsure{présch}} sous-préschémas à la fois
ouverts et fermés $P_i$\inserted{$=P\cap Q_i$} finis sur $S$. Comme ils sont étales sur $S$, leurs
}
\margin{15}

